\section{The strict permanent inequality}\label{sec:permanent}
We first prove Theorem~\ref{thm:permanent}. The two parts of the proof address different possible failures of a uniform improvement: functions close to the equality cases, and functions bounded away from them.

\begin{proof}[Proof of Theorem~\ref{thm:permanent}]
A zero function makes the assertion immediate. By homogeneity, normalize each nonzero function to have $\|f_i\|_2=1$. The resulting space of quadruples is compact. The exponent-two Carlen--Lieb--Loss inequality is strict everywhere on this space except at the single quadruple $f_i\equiv1$. Its normalization follows directly from
\[
 \operatorname{perm}(f_i(j))\le\frac{4!}{4^2}
       \prod_i\left(\sum_j f_i(j)^2\right)^{1/2}.
\]
The equality classification for nonzero nonnegative columns gives exactly the stated exceptional quadruple \cite[Theorem~1.1 and equation~(1.4)]{cll2006}.

All scalar expectations and inner products on four points use uniform measure:
$\E_jh(j)=\tfrac14\sum_jh(j)$ and $\langle g,h\rangle=\tfrac14\sum_jg(j)h(j)$.
Fix $p_*\in(4/3,2)$. Near that quadruple, the means are positive, so renormalize by the means and write $f_i=1+g_i$, with $\E_jg_i(j)=0$. Let
\[
 V=\sum_i\|g_i\|_2^2,\qquad b=\max_i\|g_i\|_\infty.
\]
For $i\ne j$, uniform sampling without replacement gives
\[
 \E_u g_i(\pi(i))g_j(\pi(j))
 =-\frac13\langle g_i,g_j\rangle.
\]
Consequently the total quadratic term is
\begin{equation}\label{eq:quadratic}
 \sum_{i<j}\E_u g_i(\pi(i))g_j(\pi(j))
 =\frac16\left(V-\Big\|\sum_i g_i\Big\|_2^2\right)
 \le\frac V6.
\end{equation}

For a law $\nu$ at total-variation distance at most $\varepsilon$ from $u$, the change in each quadratic expectation is at most
$2\varepsilon\|g_i\|_\infty\|g_j\|_\infty$. On four points, $\|g_i\|_\infty\le2\|g_i\|_2$, so the total change is at most an absolute constant times $\varepsilon V$. All linear terms are zero under $\nu$, because every coordinate marginal is uniform. Each cubic or quartic term is at most an absolute constant times $bV$ when $b\le1/2$. Indeed, bound all but two factors by $b$ and apply Cauchy--Schwarz to the remaining factors; their squared expectations use only the uniform marginals. Thus
\begin{equation}\label{eq:lhs-local}
 \E_\nu\prod_i(1+g_i(\pi(i)))
 \le1+\left(\frac16+C\varepsilon+Cb\right)V.
\end{equation}

Taylor's formula on $[1/2,3/2]$, uniformly for $p\in[p_*,2]$, gives
\[
 \E_j(1+g_i(j))^p
 =1+\frac{p(p-1)}2\|g_i\|_2^2
          +O\bigl(b\|g_i\|_2^2\bigr).
\]
Taking $p$th roots and multiplying the four factors gives
\begin{equation}\label{eq:rhs-local}
 \prod_i\|1+g_i\|_p
 =1+\frac{p-1}2 V+O(bV).
\end{equation}
All constants are uniform on the fixed interval of exponents. Since
$(p_*-1)/2>1/6$, first choose $b_*>0$ and $\varepsilon_*>0$ small enough that the error terms in \eqref{eq:lhs-local} and \eqref{eq:rhs-local} fit inside this strict gap. The desired inequality then holds for every $p\in[p_*,2]$ whenever $b<b_*$ and $\|\nu-u\|_{\TV}\le\varepsilon_*$, with exact uniform marginals.

Return to the $L^2$-normalized compact space. The condition just obtained contains a fixed open neighborhood $\mathcal U$ of its constant quadruple. On the compact complement of $\mathcal U$, the exponent-two inequality has a positive minimum gap. Both sides are continuous in the functions; they are also continuous in the exponent and in the probability masses of $\nu$. Choose $p_0\in[p_*,2)$ sufficiently close to two, then $0<\varepsilon\le\varepsilon_*$ sufficiently small, so that this gap remains positive on the complement. The local argument applies on $\mathcal U$ with these same constants. These two regions cover all normalized quadruples, proving the theorem.
\end{proof}

The exact marginal assumption can also be read off without Taylor estimates of second order. If $a(j)=\nu\{\pi(i)=j\}-1/4$ is nonzero, choose $g=a$ and set $f_i=1+t g$, with the other functions one. The derivative at $t=0$ of the left side minus the right side of \eqref{eq:permanent} is $\sum_j a(j)^2>0$. Small positive $t$ contradicts the inequality.

\begin{lemma}[Tensorization]\label{lem:tensorization}
Let $\pi_1,\ldots,\pi_b$ be independent random permutations, each with a law satisfying Theorem~\ref{thm:permanent} for the same $p_0$. If $F_i$ is a nonnegative function on $\{1,2,3,4\}^b$, then
\[
 \E\prod_{i=1}^4 F_i(\pi_1(i),\ldots,\pi_b(i))
 \le\prod_{i=1}^4
 \left(4^{-b}\sum_{x\in\{1,2,3,4\}^b}F_i(x)^{p_0}\right)^{1/p_0}.
\]
\end{lemma}
\begin{proof}
Induct on $b$. Fix the first $b-1$ coordinates and apply Theorem~\ref{thm:permanent} to the last coordinate. The four resulting functions on $b-1$ coordinates are
\[
 G_i(x)=\left(\frac14\sum_{j=1}^4 F_i(x,j)^{p_0}\right)^{1/p_0}.
\]
The induction hypothesis applies to their product, and its $p_0$th-power sums are precisely the full product-measure sums in the assertion.
\end{proof}

For later use, let $U$ be a positive semidefinite stochastic matrix of order $s$. Its entries are nonnegative and its eigenvalues lie in $[0,1]$. H\"older's inequality, applied entry by entry with powers $2-p_0$ and $p_0-1$, gives
\begin{equation}\label{eq:interpolation}
 \sum_{x,y}U(x,y)^{p_0}
 \le\left(\sum_{x,y}U(x,y)\right)^{2-p_0}
       \left(\sum_{x,y}U(x,y)^2\right)^{p_0-1}
 \le s^{2-p_0}(\Tr U)^{p_0-1}.
\end{equation}
The final inequality uses $\Tr U^2\le\Tr U$. Taking the $p_0$th root with the uniform counting normalization gives
\[
 \left(s^{-2}\sum_{x,y}U(x,y)^{p_0}\right)^{1/p_0}
 \le s^{-1}(\Tr U)^{1-1/p_0}.
\]
Section~\ref{sec:recursion} applies this estimate to four color-string transition matrices. Each contributes a factor $s^{-1}$ and a trace power $\theta=1-1/p_0$. The strict inequality $\theta<1/2$ will remain fixed throughout the proof.
